\section{Length, distance, and midpoints}

Coordinates allow us to measure geometric quantities. Once a displacement is written as a vector, its length can be computed. Once two points are written in coordinates, the distance between them can be computed. The formulas are familiar, but they should not be treated as isolated rules. They come from the Pythagorean theorem applied to coordinate differences.

\subsection*{Length of a vector}

A vector tells us how far and in which direction to move. Its length measures only the size of the displacement, not its orientation.

\begin{definitionbox}
For a vector
\[
v=(v_1,v_2)
\]
in the plane, its length is
\[
\lVert v\rVert=\sqrt{v_1^2+v_2^2}.
\]
For a vector
\[
v=(v_1,v_2,v_3)
\]
in space, its length is
\[
\lVert v\rVert=\sqrt{v_1^2+v_2^2+v_3^2}.
\]
\end{definitionbox}

The squares appear because coordinate directions are perpendicular. A displacement in the horizontal direction and a displacement in the vertical direction form the legs of a right triangle. The vector is the diagonal.

\begin{examplebox}
Let
\[
v=(3,4).
\]
Then
\[
\lVert v\rVert=\sqrt{3^2+4^2}=\sqrt{25}=5.
\]
The vector has length \(5\). The coordinates \((3,4)\) describe how the displacement is decomposed into coordinate directions; the length \(5\) describes the total size of the displacement.
\end{examplebox}

A vector and its opposite have the same length:
\[
\lVert -v\rVert=\lVert v\rVert.
\]
This is natural. Reversing a displacement changes its orientation, but not how long it is.

\subsection*{Distance between two points}

The distance between two points is the length of the vector joining them. If
\[
A=(a_1,a_2),
\qquad
B=(b_1,b_2),
\]
then
\[
\overrightarrow{AB}=(b_1-a_1,\,b_2-a_2).
\]
Therefore the distance between \(A\) and \(B\) is
\[
d(A,B)=\lVert \overrightarrow{AB}\rVert.
\]

\begin{definitionbox}
For points
\[
A=(a_1,a_2),
\qquad
B=(b_1,b_2)
\]
in the plane,
\[
d(A,B)=\sqrt{(b_1-a_1)^2+(b_2-a_2)^2}.
\]
For points in space,
\[
A=(a_1,a_2,a_3),
\qquad
B=(b_1,b_2,b_3),
\]
we have
\[
d(A,B)=\sqrt{(b_1-a_1)^2+(b_2-a_2)^2+(b_3-a_3)^2}.
\]
\end{definitionbox}

The formula is not new information. It is the length formula applied to the displacement from \(A\) to \(B\).

\begin{examplebox}
Let
\[
A=(1,-2),
\qquad
B=(5,1).
\]
Then
\[
\overrightarrow{AB}=(5-1,\,1-(-2))=(4,3).
\]
Hence
\[
d(A,B)=\lVert \overrightarrow{AB}\rVert=\sqrt{4^2+3^2}=5.
\]
\end{examplebox}

This calculation should be read in two layers. Algebraically, we subtract coordinates and take a square root. Geometrically, we first form the displacement from one point to the other, and then measure its length.

\subsection*{Midpoints}

A midpoint is the point halfway between two given points. In coordinates, it is obtained by averaging corresponding coordinates.

\begin{definitionbox}
For points
\[
A=(a_1,a_2),
\qquad
B=(b_1,b_2),
\]
the midpoint of the segment \(AB\) is
\[
M=\left(\frac{a_1+b_1}{2},\,\frac{a_2+b_2}{2}\right).
\]
In space, for
\[
A=(a_1,a_2,a_3),
\qquad
B=(b_1,b_2,b_3),
\]
the midpoint is
\[
M=\left(\frac{a_1+b_1}{2},\,\frac{a_2+b_2}{2},\,\frac{a_3+b_3}{2}\right).
\]
\end{definitionbox}

The formula is easy to remember, but again we should understand it. Each coordinate is placed halfway between the corresponding coordinates of \(A\) and \(B\).

\begin{examplebox}
Let
\[
A=(-1,4),
\qquad
B=(5,2).
\]
The midpoint is
\[
M=\left(\frac{-1+5}{2},\,\frac{4+2}{2}\right)=(2,3).
\]
The first coordinate \(2\) is halfway between \(-1\) and \(5\). The second coordinate \(3\) is halfway between \(4\) and \(2\).
\end{examplebox}

There is also a vector way to read the midpoint. Starting from \(A\), the full displacement to \(B\) is \(B-A\). Half of that displacement is \(\frac12(B-A)\). Therefore
\[
M=A+\frac12(B-A).
\]
After simplifying coordinates, this gives the same midpoint formula. This form is often more meaningful: the midpoint is reached by starting at \(A\) and moving halfway toward \(B\).

\subsection*{What is being measured?}

Length, distance, and midpoint calculations are common, but they answer different questions.

A length belongs to a vector. It asks how large a displacement is. A distance belongs to a pair of points. It asks how far apart two positions are. A midpoint is itself a point. It asks where the halfway position lies.

\begin{summarybox}
When measuring with coordinates, ask:
\begin{itemize}
\item Am I measuring a vector or comparing two points?
\item Have I first formed the correct displacement?
\item Does the result represent a number, such as a distance, or a point, such as a midpoint?
\item Can the calculation be interpreted geometrically after it is done?
\end{itemize}
\end{summarybox}

Analytic geometry becomes clearer when every formula is connected to the object it measures.